\section{Introduction}\label{sec:intro}
Fourier's law makes the heat flux proportional to the temperature gradient at the same instant. Under ordinary conditions this describes conduction in solids very well. Combined with the energy balance, the law yields a parabolic equation, and that equation has a physically uncomfortable property: a temperature disturbance applied at one point changes the temperature everywhere else at the same instant, however small the change may be.

For slow processes in bulk matter this instantaneous response cannot be distinguished from that of a slower, finite-speed carrier of heat. The difference becomes measurable, and sometimes dominant, in four situations: heating or loading that is very fast, as in femtosecond to nanosecond laser pulses; a structure that is very small or very cold, so that phonon mean free paths approach the structure size; a strongly heterogeneous medium; and a high-frequency elastic wave whose period is comparable to a thermal relaxation time.

Thermoelasticity couples the thermal field to deformation. Heating produces thermal stress through thermal expansion, and deformation changes the temperature through the dilatation term in the energy balance. The classical coupled theory of Biot \citep{Biot1956} builds on the thermomechanical stress--strain law of Duhamel and Neumann \citep{Duhamel1837} and on the early dynamic solutions of Danilovskaya \citep{Danilovskaya1950} and of Chadwick and Sneddon \citep{ChadwickSneddon1958,Chadwick1960}, and it inherits the infinite thermal signal speed. The extensions that remove that speed are called \emph{generalized} or \emph{non-classical} thermoelasticity \citep{Chandrasekharaiah1986,Chandrasekharaiah1998,HetnarskiIgnaczak1999,HetnarskiIgnaczak2000,IgnaczakOstoja2010}; they now form one of the largest strands of continuum mechanics, with a correspondingly large theoretical literature.

Two further developments widened the field. Devices exploit the coupling of deformation and temperature to further physical fields: electric polarization comes first (piezoelectricity and pyroelectricity), but magnetic induction, diffusing species, fluid-filled pores and microstructural kinematics are also used \citep{Mindlin1974Thermopiezo,Nowacki1978,Eringen1999,Iesan2004}. Independently of the heat law, the local stress--strain relation itself fails to describe dispersion at small scales, and nonlocal and gradient enrichments are used instead \citep{Eringen1972,Eringen1983,Aifantis1992,Askes2011}.

The resulting literature is large and still growing. \Cref{fig:refyears} shows the publication years of the 425 works cited here: 334 date from 1990 onward and 134 from 2020 onward, the period in which the number of named heat-conduction models for thermoelastic waves multiplied.

\begin{figure}[t]
\centering
\includegraphics[width=0.78\linewidth]{figs/fig_refyears.pdf}
\caption{Publication years of the works cited in this review, in five-year bins and colored by era. The few early bars are the foundational papers; 334 of the 425 cited works date from 1990 onward and 134 from 2020 onward, and the last bin covers only 2025--2026. The count is of cited works, not a bibliometric survey of the field.}\label{fig:refyears}
\end{figure}

Good reviews of generalized thermoelasticity already exist \citep{Chandrasekharaiah1986,Chandrasekharaiah1998,HetnarskiIgnaczak1999,HetnarskiIgnaczak2000,Shakeriaski2021,Huang2025}. What they do not supply is a treatment for a reader who wants to \emph{use} the theories on wave problems and to compare them reliably. Three problems in that literature shaped the organization of this review.
The first is terminological. The number of named heat-conduction models for thermoelastic waves has grown very large over three decades, and many of them are combinations: two-temperature fractional Lord--Shulman, three-phase-lag with memory, nonlocal Moore--Gibson--Thompson, and so on. Each paper brings its own constants and notation. \cref{sec:framework-link} introduces one heat-flux function $G(s)$ in the Laplace domain that covers every model of the three-lag-plus-memory class; Green--Lindsay and Guyer--Krumhansl do not pass through it and are treated separately, and a new combination within the class is a new choice of $G$.
The second is unreliable folklore about limits and stability. Papers repeat several statements without re-deriving them: which parameter limits reduce one model to another, which parameter ranges are stable, and how the wave branches should be labeled. These statements are re-derived here and checked numerically. One widely repeated claim, that dual-phase-lag thermoelasticity needs $\tau_T\ge\tau_q/2$ to be stable, depends on the order of the Taylor expansion that defines the model (\cref{sec:lags}).
The third is that figures taken from different papers cannot be compared. Published dispersion and response curves use different nondimensionalizations, parameter sets and branch conventions, so neither ``the effect of the fractional order'' nor ``the effect of nonlocality'' can be judged by eye. Each figure here states its own scaling (frequency in units of $\omega_*$, or lengths in units of a thermal wavelength), and the nomenclature lists the material-independent groups that control it.
The review is built on one heat-flux template that carries the classical, Lord--Shulman, Green--Naghdi, dual-phase-lag, three-phase-lag, Moore--Gibson--Thompson, fractional and memory-dependent theories, and it records which limits are regular and which are singular (Sections~\ref{sec:part_foundations}--\ref{sec:part_beyond}). For the coupled problem it derives a dispersion polynomial and compares phase speed, attenuation and thermoelastic damping, including transient, surface, guided and interface waves (\cref{sec:part_waves}). It then extends the treatment to piezoelectric, magnetic, diffusive and microstructural couplings and to nonlocal and gradient size effects (Sections~\ref{sec:part_beyond} and~\ref{sec:part_fields}). A final part covers solution methods and benchmarks, applications, experimental evidence, and a synthesis with a pitfalls checklist, a decision guide and research gaps (Sections~\ref{sec:part_methods}--\ref{sec:part_concl}).

A reader who wants to apply these theories to a wave problem should be able to choose a model and to know where it fails, and to report results in a form that other studies can compare with. That is the aim of this review.
